% Relativistic Chapter XI, §105–106: Rotation and Magnetic Field Effects on KH
% Relativistic generalization of Chandrasekhar, Chapter XI §105–106
% Conventions: see RELATIVISTIC_CONVENTIONS.md

\section{The effect of rotation on the relativistic Kelvin--Helmholtz instability}
\label{sec:rel-11-105}

We now extend the relativistic Kelvin--Helmholtz analysis of the preceding
sections to include the effect of uniform rotation about the $z$-axis with
angular velocity $\Omega$.  In the classical theory (\S\,105), the Coriolis
force modifies the perturbation equations by coupling the $x$- and
$y$-components of the velocity, and introduces the dimensionless rotation
parameter $\omega^{2}=4\Omega^{2}/gk$.  In the relativistic regime,
the Coriolis coupling must be formulated covariantly, and new
constraints arise from the requirement that all characteristic speeds
remain subluminal.

\subsection{Relativistic Coriolis coupling in a rotating frame}
\label{sec:rel-11-105a}

Consider two semi-infinite relativistic perfect fluids separated by
a horizontal interface at $z=0$.  Let the fluids have relativistic
enthalpy densities $w_{j}=(\varepsilon_{j}+p_{j})/c^{2}$ ($j=1,2$)
and stream with four-velocities whose spatial projections give
horizontal velocities $U_{1}$ and $U_{2}$ along the $x$-axis.  The
entire system rotates uniformly about the $z$-axis with angular
velocity~$\Omega$, subject to the causality requirement
\begin{equation}
\label{eq:rel-11-R1}
\boxed{r\,\Omega < c \quad\text{for all } r \text{ of interest.}}
\end{equation}

In the co-rotating frame (valid when~\eqref{eq:rel-11-R1} holds),
the linearised relativistic Euler equations for perturbations with
the $(x,y,t)$-dependence $\exp[\mathrm{i}(k_{x}x+k_{y}y+nt)]$ become
\begin{equation}
\label{eq:rel-11-R2}
\mathrm{i}\,\Lf_{j}^{2}\,w_{j}\,(n+k_{x}U_{j})\,\delta u
+ w_{j}\,(DU)\,\delta w - 2\,\Lf_{j}^{2}\,w_{j}\,\Omega\,\delta v
= -\mathrm{i}k_{x}\,\delta p,
\end{equation}
\begin{equation}
\label{eq:rel-11-R3}
\mathrm{i}\,\Lf_{j}^{2}\,w_{j}\,(n+k_{x}U_{j})\,\delta v
+ 2\,\Lf_{j}^{2}\,w_{j}\,\Omega\,\delta u
= -\mathrm{i}k_{y}\,\delta p,
\end{equation}
where $\Lf_{j}=(1-U_{j}^{2}/c^{2})^{-1/2}$ is the Lorentz factor
of fluid~$j$ and $D\equiv d/dz$.  Compared with the classical
equations~(136)--(137), each inertial term carries an additional
factor $\Lf_{j}^{2}$ reflecting the relativistic enhancement of
the effective inertia density $\Lf^{2}w$ in place of the
Newtonian~$\rho$.

Following the same algebraic steps as in the classical derivation
(eliminating the horizontal velocities and the pressure perturbation
in favour of the vertical displacement $\delta w$), one arrives at
the relativistic generalisation of equation~(143):
\begin{multline}
\label{eq:rel-11-R4}
\mathrm{i}\,\Lf^{2}\,w\,(n+k_{x}U)\!\left[
1-\frac{4\,\Lf^{4}\,\Omega^{2}}{(n+k_{x}U)^{2}}
\right]\!D(\delta w) \\
-\,\mathrm{i}\,\Lf^{2}\,w\!\left(k_{x}
+\frac{2\,\mathrm{i}\,\Lf^{2}\,k_{x}\,\Omega}{n+k_{x}U}\right)
\!(DU)\,\delta w
= -k^{2}\,\delta p.
\end{multline}
Here $k^{2}=k_{x}^{2}+k_{y}^{2}$ and we have suppressed fluid-index
subscripts for brevity.  The key relativistic modification is the
replacement $4\Omega^{2}\to 4\Lf^{4}\Omega^{2}$ in the rotation
term: the effective Coriolis parameter is enhanced by $\Lf^{2}$
because the fluid's inertia grows with the Lorentz factor.

\subsection{Two uniform relativistic fluids with rotation}
\label{sec:rel-11-105b}

In the two regions of constant $w_{j}$ and $U_{j}$,
equation~\eqref{eq:rel-11-R4} reduces to
\begin{equation}
\label{eq:rel-11-R5}
\left[1-\frac{4\,\Lf_{j}^{4}\,\Omega^{2}}{(n+k_{x}U_{j})^{2}}\right]
D^{2}(\delta w) = k^{2}\,\delta w,
\end{equation}
whose bounded solutions are $\delta w\propto e^{\pm\kappa_{j}z}$ with
\begin{equation}
\label{eq:rel-11-R6}
\kappa_{j} = \frac{k}{\sqrt{\displaystyle
1-\frac{4\,\Lf_{j}^{4}\,\Omega^{2}}{(n+k_{x}U_{j})^{2}}}}\,.
\end{equation}
Matching the continuity conditions at $z=0$ (requiring
$\delta w/(n+k_{x}U)$ continuous and integrating the governing
equation across the interface) yields the relativistic characteristic
equation
\begin{multline}
\label{eq:rel-11-R7}
\alpha_{1}^{\mathrm{rel}}\,(n+k_{x}U_{1})^{2}\!
\left[1-\frac{4\,\Lf_{1}^{4}\,\Omega^{2}}{(n+k_{x}U_{1})^{2}}
\right]^{1/2}
+\alpha_{2}^{\mathrm{rel}}\,(n+k_{x}U_{2})^{2}\!
\left[1-\frac{4\,\Lf_{2}^{4}\,\Omega^{2}}{(n+k_{x}U_{2})^{2}}
\right]^{1/2} \\
= gk\bigl(\alpha_{1}^{\mathrm{rel}}-\alpha_{2}^{\mathrm{rel}}\bigr)
+ \frac{k^{3}T}{w_{1}+w_{2}}\,,
\end{multline}
where the \emph{relativistic density fractions} are
\begin{equation}
\label{eq:rel-11-R8}
\alpha_{j}^{\mathrm{rel}}
= \frac{\Lf_{j}^{2}\,w_{j}}{\Lf_{1}^{2}\,w_{1}+\Lf_{2}^{2}\,w_{2}}\,,
\qquad j=1,2,
\end{equation}
with $\alpha_{1}^{\mathrm{rel}}+\alpha_{2}^{\mathrm{rel}}=1$.

In the non-relativistic limit $U_{j}/c\to 0$ (so $\Lf_{j}\to 1$ and
$w_{j}\to\rho_{j}$), the characteristic equation~\eqref{eq:rel-11-R7}
reduces exactly to the classical equation~(153).

\subsection{Characteristic equation analysis}
\label{sec:rel-11-105c}

Setting $T=0$, restricting to perturbations along the streaming
direction ($k_{x}=k$), and introducing the phase velocity $c_{p}=-n/k$,
equation~\eqref{eq:rel-11-R7} becomes
\begin{multline}
\label{eq:rel-11-R9}
\alpha_{1}^{\mathrm{rel}}\,(U_{1}-c_{p})^{2}
\left[1-\frac{4\,\Lf_{1}^{4}\,\Omega^{2}}{k^{2}(U_{1}-c_{p})^{2}}
\right]^{1/2}\\
+\,\alpha_{2}^{\mathrm{rel}}\,(U_{2}-c_{p})^{2}
\left[1-\frac{4\,\Lf_{2}^{4}\,\Omega^{2}}{k^{2}(U_{2}-c_{p})^{2}}
\right]^{1/2}
= \frac{g}{k}\bigl(\alpha_{1}^{\mathrm{rel}}-\alpha_{2}^{\mathrm{rel}}\bigr).
\end{multline}
Introducing the dimensionless variables
$\xi=(U_{1}-c_{p})\sqrt{k/g}$, $\eta=(U_{2}-c_{p})\sqrt{k/g}$, and
the \emph{relativistic rotation parameter}
\begin{equation}
\label{eq:rel-11-R10}
\omega_{j}^{2} = \frac{4\,\Lf_{j}^{4}\,\Omega^{2}}{gk}\,,
\qquad j=1,2,
\end{equation}
equation~\eqref{eq:rel-11-R9} takes the same structural form as the
classical equation~(158), but with $\omega\to\omega_{j}$ depending on
the fluid label (the two fluids have different effective rotation
parameters when their streaming velocities differ).

The graphical analysis of the classical theory (\S\,105\,(b))---the
$P$-, $S_{1}$-, $S_{2}$-, and $C$-branches in the
$(\xi,\eta)$-plane---carries over verbatim with $\omega$ replaced by
the fluid-dependent $\omega_{j}$.  The singular points become
$(\xi_{0},\omega_{2})$ and $(\omega_{1},\eta_{0})$, and the principal
result is:

\smallskip
\noindent\textit{For a system of two relativistic fluids partaking in
rotation, the Kelvin--Helmholtz instability is suppressed so long as}
\begin{equation}
\label{eq:rel-11-R11}
\boxed{
|V|\,\sqrt{\frac{k}{g}}
\leqslant
\left\{\tfrac{1}{2}\,\omega_{2}^{2}
+\tfrac{1}{2}\!\left[\omega_{2}^{4}
+4\!\left(\frac{\alpha_{1}^{\mathrm{rel}}
-\alpha_{2}^{\mathrm{rel}}}{\alpha_{2}^{\mathrm{rel}}}
\right)^{\!2}\right]^{1/2}\right\}^{1/2}\!\!-\omega_{1}\,,}
\end{equation}
\textit{where $V=U_{1}-U_{2}$ is the relative streaming velocity.}

\subsubsection*{Causality constraint}

The phase velocity of any mode must satisfy $|c_{p}|<c$.  From
equation~\eqref{eq:rel-11-R9}, the maximum phase velocity of the
stable modes is bounded above by $\max(U_{1},U_{2})+\eta_{0}\sqrt{g/k}$.
Causality therefore imposes an additional global condition:
\begin{equation}
\label{eq:rel-11-R12}
\boxed{\max(|U_{1}|,\,|U_{2}|)
+ \eta_{0}\,\sqrt{g/k}\;<\;c\,.}
\end{equation}
When this bound is saturated, the perturbation can no longer be treated
as a subluminal hydrodynamic mode; relativistic compressibility and
radiation effects would need to be included.


%% ================================================================
\section{The effect of a horizontal magnetic field on the relativistic
  Kelvin--Helmholtz instability}
\label{sec:rel-11-106}

We now consider the relativistic analogue of the magnetic
stabilisation of the Kelvin--Helmholtz instability (\S\,106).  A
uniform magnetic field, when aligned with the direction of streaming,
introduces the Alfv\'en speed as a restoring agent; the competition
between the destabilising velocity shear and the stabilising magnetic
tension determines the onset of instability.

\subsection{Relativistic Alfv\'en speed and magnetic tension}
\label{sec:rel-11-106a}

In the framework of relativistic ideal MHD (cf.\ the conventions in
\textsc{Relativistic\_Conventions.md}), the energy-momentum tensor of
a magnetised perfect fluid is
\begin{equation}
\label{eq:rel-11-R13}
T^{\mu\nu}_{\mathrm{MHD}}
= \left(w+\frac{b^{2}}{c^{2}}\right)u^{\mu}u^{\nu}
+ \left(p+\frac{b^{2}}{2}\right)g^{\mu\nu}
- b^{\mu}b^{\nu},
\end{equation}
where $b^{\mu}$ is the magnetic four-vector in the fluid rest frame
($b^{\mu}u_{\mu}=0$) and $b^{2}=b^{\mu}b_{\mu}$.  The relativistic
Alfv\'en speed is
\begin{equation}
\label{eq:rel-11-R14}
\vA^{2}
= \frac{b^{2}}{4\pi\,w + b^{2}}\,c^{2}
= \frac{v_{A,\mathrm{class}}^{2}}
{1+v_{A,\mathrm{class}}^{2}/c^{2}}\,,
\end{equation}
where $v_{A,\mathrm{class}}^{2}=\mu H^{2}/(4\pi\rho)$ is the
classical Alfv\'en speed and $w=(\varepsilon+p)/c^{2}$ is the
relativistic enthalpy density.  By construction,
\begin{equation}
\label{eq:rel-11-R15}
\boxed{\vA < c\,,}
\end{equation}
which is the fundamental causality bound for Alfv\'en waves.

\subsection{Two uniform magnetised fluids: parallel field}
\label{sec:rel-11-106b}

Consider two semi-infinite relativistic fluids ($z<0$ and $z>0$) with
enthalpy densities $w_{1}$, $w_{2}$, streaming velocities $U_{1}$,
$U_{2}$ along $x$, and a uniform magnetic field $\mathbf{H}=H\,\hat{x}$
also aligned with the streaming direction.  The linearised relativistic
MHD perturbation equations, for modes with the $(x,y,t)$-dependence
$\exp[\mathrm{i}(k_{x}x+k_{y}y+nt)]$, yield bounded solutions
analogous to the classical equations~(199)--(200).

Matching boundary conditions at $z=0$, we obtain the relativistic
characteristic equation
\begin{equation}
\label{eq:rel-11-R16}
\Lf_{1}^{2}\,w_{1}\,(n+k_{x}U_{1})^{2}
+\Lf_{2}^{2}\,w_{2}\,(n+k_{x}U_{2})^{2}
= gk\bigl(\Lf_{1}^{2}\,w_{1}-\Lf_{2}^{2}\,w_{2}\bigr)
+ k_{x}^{2}\!\left(\frac{b_{1}^{2}}{4\pi}+\frac{b_{2}^{2}}{4\pi}\right),
\end{equation}
where $b_{j}^{2}$ is the comoving magnetic-pressure parameter of
fluid~$j$.  The roots are
\begin{multline}
\label{eq:rel-11-R17}
n = -k_{x}\bigl(\alpha_{1}^{\mathrm{rel}}U_{1}
+\alpha_{2}^{\mathrm{rel}}U_{2}\bigr) \\
\pm\left[gk\bigl(\alpha_{1}^{\mathrm{rel}}-\alpha_{2}^{\mathrm{rel}}\bigr)
+ k_{x}^{2}\,\frac{b_{1}^{2}+b_{2}^{2}}
{4\pi(\Lf_{1}^{2}w_{1}+\Lf_{2}^{2}w_{2})}
- k_{x}^{2}\,\alpha_{1}^{\mathrm{rel}}\,\alpha_{2}^{\mathrm{rel}}\,
(U_{1}-U_{2})^{2}\right]^{1/2}\!.
\end{multline}

The instability is suppressed when the discriminant is non-negative.
For perturbations in the streaming direction ($k_{x}=k$) with
$\alpha_{1}^{\mathrm{rel}}\geqslant\alpha_{2}^{\mathrm{rel}}$, the
critical condition is:

\smallskip
\noindent\textit{A uniform magnetic field parallel to the direction of
streaming suppresses the relativistic Kelvin--Helmholtz instability if
and only if}
\begin{equation}
\label{eq:rel-11-R18}
\boxed{
\alpha_{1}^{\mathrm{rel}}\,\alpha_{2}^{\mathrm{rel}}\,(U_{1}-U_{2})^{2}
\;\leqslant\;
\frac{b_{1}^{2}+b_{2}^{2}}
{4\pi\bigl(\Lf_{1}^{2}\,w_{1}+\Lf_{2}^{2}\,w_{2}\bigr)}\,,}
\end{equation}
\textit{i.e.\ the relative streaming velocity must not exceed the
root-mean-square relativistic Alfv\'en speed in the two media.}

Noting that $b_{j}^{2}/(4\pi w_{j})=\vA_{,j}^{2}/c^{2}\cdot w_{j}$
and using~\eqref{eq:rel-11-R14}, the stability criterion can be
rewritten as
\begin{equation}
\label{eq:rel-11-R19}
(U_{1}-U_{2})^{2}
\leqslant
\frac{\Lf_{1}^{2}\,w_{1}\,\vA_{,1}^{2}
+\Lf_{2}^{2}\,w_{2}\,\vA_{,2}^{2}}
{\Lf_{1}^{2}\,w_{1}\,\alpha_{2}^{\mathrm{rel}}
\cdot c^{2}\cdot\alpha_{1}^{\mathrm{rel}}}\,,
\end{equation}
which manifestly reduces to the classical criterion~(205) in the
limit $c\to\infty$.

\subsubsection*{Effective surface tension}

Comparing equation~\eqref{eq:rel-11-R16} with the field-free
relativistic KH equation, the magnetic field acts as an effective
surface tension of amount
\begin{equation}
\label{eq:rel-11-R20}
T_{\mathrm{eff}}^{\mathrm{rel}}
= \frac{b_{1}^{2}+b_{2}^{2}}{4\pi}\,\frac{k_{x}^{2}}{k^{2}}\,,
\end{equation}
which is the relativistic counterpart of equation~(203).  The field
stabilises only perturbations with a component along the field
direction ($k_{x}\neq 0$); purely transverse modes ($k_{x}=0$) are
unaffected.

\subsection{Transverse magnetic field: anisotropic stabilisation}
\label{sec:rel-11-106c}

When the impressed magnetic field $\mathbf{H}=H\,\hat{y}$ is
transverse to the streaming direction, the relativistic MHD
perturbation equations reduce, after lengthy but straightforward
algebra, to a form that is identical to the field-free relativistic
KH equation for perturbations propagating along the streaming
direction.  This is the relativistic generalisation of
equation~(210):

\smallskip
\noindent\textit{The development of the relativistic
Kelvin--Helmholtz instability in the direction of streaming is
uninfluenced by a magnetic field transverse to that direction.}

\smallskip
The physical explanation is the same as in the classical case: the
magnetic tension from a transverse field acts only on perturbations
that bend the field lines, i.e.\ those with a wavevector component
along $\hat{y}$.  Perturbations propagating purely along $\hat{x}$
(the streaming direction) leave the transverse field lines unperturbed
and receive no restoring force.

For oblique perturbations (both $k_{x}\neq 0$ and $k_{y}\neq 0$),
the stabilisation is anisotropic: the magnetic tension provides a
restoring force proportional to $k_{y}^{2}$ from the transverse
field, which combines with the shear-driven instability proportional
to $k_{x}^{2}$.  The instability is then confined to wavevectors
within a cone around $\hat{x}$ whose half-angle shrinks as the
magnetic field strength increases.

\subsection{Causality: combined Alfv\'en and KH mode speeds}
\label{sec:rel-11-106d}

The relativistic treatment imposes a non-trivial constraint that has
no classical counterpart.  The phase velocity of the coupled
Alfv\'en--KH modes must satisfy
\begin{equation}
\label{eq:rel-11-R21}
\boxed{|c_{p}| = \left|\frac{n}{k}\right| < c\,.}
\end{equation}
From equation~\eqref{eq:rel-11-R17}, the maximum phase speed of a
stable mode is
\begin{equation}
\label{eq:rel-11-R22}
|c_{p}|_{\max}
\leqslant
\alpha_{1}^{\mathrm{rel}}|U_{1}|
+\alpha_{2}^{\mathrm{rel}}|U_{2}|
+\left[\frac{g}{k}\bigl(\alpha_{1}^{\mathrm{rel}}
-\alpha_{2}^{\mathrm{rel}}\bigr)
+\vA^{2}_{\mathrm{rms}}\right]^{1/2}\!,
\end{equation}
where $\vA^{2}_{\mathrm{rms}}$ is the enthalpy-weighted
root-mean-square relativistic Alfv\'en speed.  Since each
$|U_{j}|<c$ and each $\vA_{,j}<c$ by construction, the causality
bound~\eqref{eq:rel-11-R21} is automatically satisfied whenever
the individual fluid velocities and Alfv\'en speeds are subluminal.
In the formal limit where either $|U_{j}|\to c$ or $\vA_{,j}\to c$,
the incompressible treatment breaks down and a fully compressible
relativistic MHD analysis is required.

More precisely, the fast magnetosonic speed
\begin{equation}
\label{eq:rel-11-R23}
v_{f}^{2} = \cs^{2}+\vA^{2}-\frac{\cs^{2}\,\vA^{2}}{c^{2}}
\end{equation}
provides the upper bound on any MHD mode; the algebraic form of
equation~\eqref{eq:rel-11-R23} guarantees $v_{f}<c$ whenever
$\cs<c$ and $\vA<c$.

%% ================================================================
\section*{Bibliographical Notes}
\label{sec:rel-11-105-106-bib}

The classical analysis of rotation effects on the Kelvin--Helmholtz
instability is due to:

\begin{enumerate}

\item \textsc{S.\ Chandrasekhar}, `The effect of rotation on the
  development of the Kelvin--Helmholtz instability', unpublished.

\item \textsc{H.\ Solberg}, `Integrationen der
  atmosph\"arischen St\"orungsgleichungen, I.\ Wellenbewegungen in
  rotierenden, inkompressiblen Fl\"ussigkeitsschichten',
  \emph{Geofysiske Publikasjoner}, \textbf{9}, 1--120 (1928).
\end{enumerate}

\medskip\noindent
The classical hydromagnetic Kelvin--Helmholtz instability was treated by:

\begin{enumerate}
\setcounter{enumi}{2}
\item \textsc{D.\ H.\ Michael}, `The stability of a combined current
  and vortex sheet in a perfectly conducting fluid',
  \emph{Proc.\ Camb.\ Phil.\ Soc.}, \textbf{51}, 528--32 (1955).

\item \textsc{T.\ G.\ Northrop}, `Helmholtz instability of a plasma',
  \emph{Physical Rev.}, \textbf{103}, 1150--4 (1956).
\end{enumerate}

\medskip\noindent
The relativistic extension of the Kelvin--Helmholtz instability,
particularly in the context of astrophysical jets, has been developed in:

\begin{enumerate}
\setcounter{enumi}{4}
\item \textsc{A.\ Ferrari}, \textsc{E.\ Trussoni}, and
  \textsc{L.\ Zaninetti}, `Relativistic Kelvin--Helmholtz instabilities
  in radiative jets', \emph{Astron.\ Astrophys.}, \textbf{64}, 43--52
  (1978).

\item \textsc{P.\ E.\ Hardee}, `On the configuration and propagation of
  jets in extragalactic radio sources',
  \emph{Astrophys.\ J.}, \textbf{277}, 106--12 (1984).

\item \textsc{M.\ Perucho}, \textsc{J.\ M.\ Mart\'{\i}},
  \textsc{M.\ Hanasz}, `Stability of hydrodynamical relativistic
  planar jets.\ I.\ Linear evolution and saturation of Kelvin--Helmholtz
  modes', \emph{Astron.\ Astrophys.}, \textbf{427}, 415--29 (2004).

\item \textsc{M.\ Perucho}, \textsc{J.\ M.\ Mart\'{\i}},
  \textsc{J.\ M.\ Cela}, \textsc{M.\ Hanasz}, \textsc{R.\ de la Cruz},
  and \textsc{F.\ Rubio}, `Stability of three-dimensional relativistic
  jets', \emph{Astron.\ Astrophys.}, \textbf{519}, A41 (2010).
\end{enumerate}

\medskip\noindent
The role of magnetic fields in stabilising relativistic jets against KH
modes:

\begin{enumerate}
\setcounter{enumi}{8}
\item \textsc{I.\ Mizuno}, \textsc{P.\ E.\ Hardee}, and
  \textsc{K.-I.\ Nishikawa}, `Spatial growth of the current-driven and
  Kelvin--Helmholtz instability in relativistic jets',
  \emph{Astrophys.\ J.}, \textbf{734}, 19 (2011).

\item \textsc{E.\ M.\ de Gouveia Dal Pino}, `Magnetic fields in
  astrophysical jets', \emph{Adv.\ Space Res.}, \textbf{35},
  908--17 (2005).
\end{enumerate}

\medskip\noindent
The relativistic Alfv\'en speed and its role as a causality bound in
MHD:

\begin{enumerate}
\setcounter{enumi}{10}
\item \textsc{A.\ Lichnerowicz}, \emph{Relativistic Hydrodynamics and
  Magnetohydrodynamics}, W.\ A.\ Benjamin, New York, 1967.

\item \textsc{L.\ Rezzolla} and \textsc{O.\ Zanotti},
  \emph{Relativistic Hydrodynamics}, Oxford University Press, Oxford,
  2013.
\end{enumerate}

\medskip\noindent
Causality constraints on relativistic MHD perturbations and their
implications for jet stability:

\begin{enumerate}
\setcounter{enumi}{12}
\item \textsc{S.\ S.\ Komissarov}, `On the properties of Alfv\'en
  waves in relativistic magnetohydrodynamics',
  \emph{Phys.\ Lett.\ A}, \textbf{232}, 435--42 (1997).

\end{enumerate}

% ===================================================================
\section*{Astrophysical application: magnetic-field stabilisation of relativistic jets}
\label{sec:rel-11-105-106-app}

The magnetic stabilisation criterion~\eqref{eq:rel-11-R18} is of
central importance for understanding the long-term survival of
magnetised relativistic jets.  Observations and simulations indicate
that AGN jets carry dynamically significant magnetic fields, with
magnetisation parameters $\sigma = B^{2}/(4\pi w c^{2})$ ranging
from $\sigma \sim 0.01$ (kinetically dominated) to $\sigma \sim 1$
or greater near the jet base (magnetically dominated)
\cite{MizunoHardeeNishikawa2007, MizunoEtAl2011, BodoEtAl2013}.

\paragraph{Critical magnetisation.}
From~\eqref{eq:rel-11-R18}, the critical magnetisation for complete
KH stabilisation of a jet with Lorentz factor $\Gamma$ and ambient
density ratio $\eta = \hat{\rho}_{\mathrm{amb}}/\hat{\rho}_{\mathrm{jet}}$
is
\begin{equation}
  \sigma_{\mathrm{crit}}
  = \frac{\Gamma^{2}\,\eta}{(\Gamma^{2}+\eta)^{2}}\,\beta_{\mathrm{jet}}^{2}\,.
\end{equation}
For a light jet ($\eta \ll 1$) at high $\Gamma$, this reduces to
$\sigma_{\mathrm{crit}} \approx \eta/\Gamma^{2}$, showing that only
a very weak magnetic field is needed to stabilise a highly relativistic,
underdense jet.  For a matched jet ($\eta \sim 1$), the critical
magnetisation is of order $\sigma_{\mathrm{crit}} \sim 1/\Gamma^{2}$,
still achievable for moderate $\Gamma$.

\paragraph{Implications for jet morphology.}
The stabilisation by a parallel magnetic field explains several observed
features of AGN jets:
\begin{enumerate}
\item Parsec-scale jets, which are highly magnetised and show smooth,
  well-collimated structures, are consistent with
  $\sigma > \sigma_{\mathrm{crit}}$ \cite{MizunoEtAl2011}.
\item Kiloparsec-scale jets, where the magnetic field has decayed through
  expansion, show knotty, disrupted morphologies consistent with the
  onset of KH instability as $\sigma$ drops below the critical value.
\item The transition from FR\,I (edge-brightened, disrupted) to FR\,II
  (well-collimated, terminated by hotspots) jet morphology can be
  understood in terms of the magnetisation threshold: FR\,II jets
  maintain $\sigma > \sigma_{\mathrm{crit}}$ over their full length
  \cite{BodoEtAl2013}.
\end{enumerate}

Figure~\ref{fig:magnetized-jet-stabilization} shows the critical
magnetisation as a function of $\Gamma$ for different density ratios,
and the KH growth rate as a function of $\sigma$ for representative
jet Lorentz factors.

\begin{figure}[htbp]
\centering
\includegraphics[width=0.95\textwidth]{fig_magnetized_jet_stabilization.pdf}
\caption{%
\textbf{Left:} Critical magnetisation $\sigma_{\mathrm{crit}}$ for
complete KH stabilisation vs.\ jet Lorentz factor, for density ratios
$\eta = 0.01$--$1.0$.
\textbf{Right:} Normalised KH growth rate vs.\ magnetisation parameter
$\sigma$ for $\Gamma = 3$, 10, and 30 (with $\eta = 0.1$).  The
vertical dashed lines mark the critical $\sigma$ for each $\Gamma$.
}
\label{fig:magnetized-jet-stabilization}
\end{figure}

% ===================================================================
\section{Deep Research: instabilities at recollimation shocks in MHD jets}
\label{sec:rel-11-105-106-recoll}

Relativistic jets propagating through declining-pressure environments
(e.g.\ the hot gas halo of the host galaxy) undergo periodic
\emph{recollimation}: the jet over-expands, then recollimates under
external pressure, producing stationary oblique shocks at the
reconfinement points \cite{GourgouliatosKomissarov2018, Matsumoto2017}.
Recent relativistic MHD simulations by \cite{MukherjeeBodo2025}
demonstrate that these recollimation shocks profoundly alter the
local KH stability of the jet by modifying both the flow velocity
and the magnetic field structure.

\subsection{Post-shock conditions}
\label{sec:rel-11-105-106-recoll-post}

At a recollimation shock with compression ratio $\mathcal{C} \sim 3$
(typical for moderately relativistic jets), the post-shock flow
parameters are determined by the relativistic MHD Rankine--Hugoniot
conditions:
\begin{enumerate}
\item \textbf{Deceleration:} The post-shock Lorentz factor is
  $\Gamma_{\mathrm{post}} \approx \Gamma_{\mathrm{pre}}/\mathcal{C}$,
  reducing the $\Gamma^{-2}$ time-dilation stabilisation.
\item \textbf{Density compression:} $\hat{\rho}_{\mathrm{post}} \approx
  \mathcal{C}\,\hat{\rho}_{\mathrm{pre}}$, increasing the density
  contrast with the ambient.
\item \textbf{Magnetic amplification:} For a predominantly toroidal
  field, $B_{\mathrm{post}} \approx \mathcal{C}\,B_{\mathrm{pre}}$,
  so the magnetisation parameter jumps:
  $\sigma_{\mathrm{post}} \approx \mathcal{C}^{2}\,\sigma_{\mathrm{pre}}$.
\end{enumerate}

The competition between effects~(1)--(2) (which promote instability)
and effect~(3) (which stabilises) determines whether the post-shock
region is KH-unstable.

\subsection{Post-shock KH growth rate}
\label{sec:rel-11-105-106-recoll-growth}

Immediately downstream of the recollimation shock, the magnetic field
is strongly compressed and the jet is stabilised
($\sigma_{\mathrm{post}} > \sigma_{\mathrm{crit}}$ from
equation~\eqref{eq:rel-11-R18}).  However, as the jet re-expands
downstream, the field decays as $B \propto r^{-1}$ (toroidal field
in conical expansion), so $\sigma \propto r^{-2}$.

The critical distance at which the magnetisation drops below the
stabilisation threshold defines the \emph{KH onset distance}
downstream of the shock:
\begin{equation}\label{eq:rel-recoll-1}
  z_{\mathrm{onset}} \approx R_{\mathrm{jet}}\,
  \sqrt{\frac{\sigma_{\mathrm{post}}}{\sigma_{\mathrm{crit}}}}
  \cdot \frac{1}{\tan\theta_{\mathrm{expand}}}\,,
\end{equation}
where $\theta_{\mathrm{expand}}$ is the post-shock re-expansion angle.
For typical parameters ($\sigma_{\mathrm{pre}} = 0.1$,
$\mathcal{C} = 3$, $\theta_{\mathrm{expand}} \sim 3^{\circ}$),
the onset occurs at $z_{\mathrm{onset}} \sim 5$--$10\,R_{\mathrm{jet}}$
downstream.

Beyond this distance, the hydrodynamic KH growth rate applies, but
with the modified post-shock flow parameters:
\begin{equation}\label{eq:rel-recoll-2}
  \sigma_{\mathrm{KH}}(z)
  \sim \frac{k\,c_s\,\sqrt{\eta_{\mathrm{post}}}}
  {\Gamma_{\mathrm{post}}^{2}}\,
  \left(1 - \frac{\sigma_{B}(z)}{\sigma_{\mathrm{crit}}(z)}\right),
\end{equation}
which increases with distance as $\sigma_{B}(z)$ decays.

The simulations of \cite{MukherjeeBodo2025} confirm this analytical
picture: the jet develops KH-driven knotty structures at distances
$\sim 5$--$15\,R_{\mathrm{jet}}$ downstream of each recollimation
shock, with the most prominent knots appearing where the local
magnetisation first drops below the critical value.  This provides
a natural explanation for the quasi-periodic knotty structures
observed in many AGN jets.

Figure~\ref{fig:recollimation-kh} shows the post-recollimation KH
growth rate as a function of downstream distance for magnetised
relativistic jets, together with the magnetic stabilisation boundary.

\begin{figure}[htbp]
\centering
\includegraphics[width=0.95\textwidth]{fig_recollimation_kh.pdf}
\caption{%
\textbf{Left:} KH growth rate downstream of a recollimation shock
($\Gamma_{\mathrm{pre}} = 10$, $\sigma_{\mathrm{pre}} = 0.1$,
compression ratio $\mathcal{C} = 3$).  Solid red: MHD (with decaying
field); dashed blue: hydrodynamic (no field).  The magnetisation
$\sigma_{B}(z)$ and critical magnetisation $\sigma_{\mathrm{crit}}(z)$
are shown on the right axis.  The vertical dotted line marks the KH
onset distance.
\textbf{Right:} Critical magnetisation boundary downstream of the
shock for different pre-shock Lorentz factors
($\Gamma_{\mathrm{pre}} = 5$, 10, 20); the dashed black curve
shows the actual $\sigma_{B}(z)$ trajectory.  The star marks the
simulation result of \cite{MukherjeeBodo2025}.
}
\label{fig:recollimation-kh}
\end{figure}

%% =====================================================================
%%  Gavassino integration: Charge separation and superflow stability
%%  Papers 44 (2508.08936) and 41 (2507.19985) — brief cross-references
%% =====================================================================

\begin{relcorrection}
\textbf{Charge-neutrality breakdown and superflow analogy.}
The Kelvin--Helmholtz analyses above assume MHD charge neutrality in
the background state.  As discussed in
\S\ref{sec:rel-10-charge-separation}, this assumption can fail for
relativistic plasmas in strong gravitational fields
\cite{GavassinoChargeSeparation2025}.  Additionally, the
thermodynamic stability framework for superflows developed by
Gavassino~\cite{GavassinoSuperflows2025} (see
\S\ref{sec:rel-12-superflow-stability}) provides a complementary
perspective on the KH instability: the onset of instability at a
shear interface can be viewed as the point where the effective grand
potential ceases to be minimised by the sheared equilibrium state,
connecting the KH criterion to the broader principle that
long-lived flows must be thermodynamic minima.
\end{relcorrection}


%% ====================================================================
\subsection{Acausality-driven instabilities in Israel--Stewart and validation of BDNK}
\label{sec:rel-11-IS-acausality-KH}
%% ====================================================================

The relativistic Kelvin--Helmholtz instability is particularly
susceptible to the acausality-driven pathologies of Israel--Stewart
theory identified by Gavassino et~al.\
\citep{GavassinoAcausalityIS2025}.  In the KH configuration,
the velocity shear layer generates large viscous stresses, and the
streaming velocities $U_{1}$, $U_{2}$ are themselves relativistic.
Both factors conspire to push the IS characteristic speed
$w$ (eq.~\ref{eq:IS-info-speed-RT}) above unity and the product $vw$
above the instability threshold $vw \ge 1$, where $v$ is the local
fluid speed in the computational frame.

When this ``ugly'' regime is reached, the IS equations suffer from
reversed dissipation: perturbations at the shear interface grow not
because of the physical KH mechanism but because of backward-propagating
signals that amplify noise exponentially
\citep{GavassinoAcausalityIS2025,GavassinoCausalityHIC2025}.  In
heavy-ion collision simulations---where IS theory is the standard
tool---Gavassino et~al.\ demonstrated that significant fractions
of the computational domain can enter the ``bad'' or ``ugly'' regime
during the early high-gradient evolution phase.

For astrophysical jet simulations of KH instability, the situation
is even more severe: the high Lorentz factors ($\Gamma \sim 5$--$20$)
and large pressure jumps at the jet boundary make it nearly impossible
to avoid acausal regions within IS theory.  The resulting spurious
instabilities would contaminate the physical KH growth rates and
saturated mixing-layer widths.

The BDNK framework avoids this pathology entirely.  Because the
BDNK characteristic speeds depend only on the transport coefficients
$(\eta, \zeta, \kappa)$ and thermodynamic state---not on the viscous
stresses or the flow velocity---they remain subluminal throughout the
nonlinear KH evolution.  This provides strong validation for the
exclusive use of BDNK dissipation in all relativistic KH analyses
presented in this chapter.
